\clearpage\begingroup\SourceLeaf{68b}\HistoricalHead{63}{R}{BERECHNUNG MAGNETISCHER KRÄFTE.}
\begin{center}
\rule{.25\textwidth}{.4pt}\\
\textbf{[Aus Handbuch 19, Be.]}\\[-.3ex]
\rule{.25\textwidth}{.4pt}
\end{center}
\begin{center}[6.]\end{center}
\noindent{\small [S. 199]}
\begin{center}\textbf{Aufgabe.}\end{center}
\indent Fünf Commutationen des galvanischen Stromes so anzuordnen, dass zwei
in der Kette befindliche Nadeln nachher eben so grosse Schwingungen haben,
wie vorher.

Es sein \(\dfrac{180^0}{n},\ \dfrac{180^0}{n'}\) die Schwingungszeiten der beiden Nadeln,
\[
T-t,\ T-u,\ T,\ T+u',\ T+t'
\]
\clearpage\SourceLeaf{69a}\HistoricalHead{64}{L}{PHYSIK. NACHLASS.}
\noindent die Augenblicke der fünf Commutationen. Die vorgeschriebene Bedingung
erfordert, dass
\[
\begin{array}{r@{}c@{}l}
\cos nt-\cos nu&{}+1{}&-\cos nu'+\cos nt'=0,\\
\sin nt-\sin nu&&+\sin nu'-\sin nt'=0,\\
\cos n't-\cos n'u&{}+1{}&-\cos n'u'+\cos n't'=0,\\
\sin n't-\sin n'u&&+\sin n'u'-\sin n't'=0.
\end{array}
\]
\noindent Der zweiten und vierten Gleichung wird Genüge geleistet, indem man \(u'=u,\)
\(t'=t\) setzt. Wir haben demnach
\[
\begin{aligned}
\cos nu-\cos nt&=\tfrac12,\\
\cos n'u-\cos n't&=\tfrac12,
\end{aligned}
\]
\noindent woraus sich \(t\) und \(u\) leicht durch Versuche bestimmen lassen. So findet sich
für die Schwingungsdauern \(20^{\prime\prime}6\) und \(43^{\prime\prime}1\)
\[
t=21^{\prime\prime}00,\quad u\mathpunct{.}=13^{\prime\prime}72.
\]
\noindent Die Zeiten der Commutationen folgen also so aufeinander
\[
0;\ 7^{\prime\prime}28;\ 21^{\prime\prime}00;\ 34^{\prime\prime}72;\ 42^{\prime\prime}00.
\]
\clearpage\endgroup
